\section{Questao 5 - Minimax e Poda Alpha-Beta}

\subsection{Conceitos pedidos no enunciado}
\begin{itemize}[leftmargin=1.5em]
  \item Objetivo do Minimax: escolher a acao que maximiza o resultado de MAX assumindo MIN otimo.
  \item Diferenca MAX/MIN: nos MAX escolhem maior valor; nos MIN escolhem menor valor.
  \item Utilidade: valor terminal (folha) que mede o resultado final do jogo.
  \item Propagacao: os valores sobem das folhas para a raiz alternando regras MAX/MIN.
  \item Adversario perfeito: MIN sempre joga da melhor forma para reduzir o ganho de MAX.
  \item Poda Alpha-Beta: elimina ramos que nao podem mais alterar a decisao final da raiz.
\end{itemize}

\subsection{Dados da arvore}
Folhas (esquerda para direita):
\[
[3,5,6,9,1,2,0,-1,7,4,5,6]
\]
Raiz $A$ e do tipo MAX; filhos $B$, $C$, $D$ sao MIN; cada um possui 4 folhas.

\subsection{Minimax passo a passo}
\begin{table}[H]
\centering
\begin{tabular}{c c c c}
\toprule
No & Tipo & Folhas/filhos avaliados & Valor \\
\midrule
$B$ & MIN & $[3,5,6,9]$ & 3 \\
$C$ & MIN & $[1,2,0,-1]$ & -1 \\
$D$ & MIN & $[7,4,5,6]$ & 4 \\
$A$ & MAX & $[B=3,C=-1,D=4]$ & 4 \\
\bottomrule
\end{tabular}
\caption{Propagacao de valores do Minimax}
\end{table}

\textbf{Decisao de MAX na raiz:} escolher $D$ (valor 4).\\
\textbf{Ordem de expansao usada:} esquerda $\rightarrow$ direita (B, depois C, depois D; e folhas de cada no tambem da esquerda para direita).\\
\textbf{Caminho escolhido:} $A \rightarrow D$.\\
\textbf{Nos visitados:} 16 (4 internos + 12 folhas).

\subsection{Arvore parcialmente preenchida (visao estrutural)}
\begin{figure}[H]
\centering
\begin{forest}
for tree={
  draw,
  rounded corners,
  align=center,
  parent anchor=south,
  child anchor=north,
  l sep=10mm,
  s sep=8mm
}
[A\\MAX\\4
  [B\\MIN\\3
    [3]
    [5]
    [6]
    [9]
  ]
  [C\\MIN\\-1
    [1]
    [2]
    [0]
    [-1]
  ]
  [D\\MIN\\4
    [7]
    [4]
    [5]
    [6]
  ]
]
\end{forest}
\caption{Arvore da questao com valores Minimax}
\end{figure}

\subsection{Alpha-Beta (mesma ordem do enunciado)}
\begingroup
\small
\begin{longtable}{c c c c c c}
\toprule
Passo & No MIN & Folha & $\alpha$ & $\beta$ & Poda \\
\midrule
1 & B & 3 & $-\infty$ & 3 & nao \\
2 & B & 5 & $-\infty$ & 3 & nao \\
3 & B & 6 & $-\infty$ & 3 & nao \\
4 & B & 9 & $-\infty$ & 3 & nao \\
5 & C & 1 & 3 & 1 & sim \\
6 & D & 7 & 3 & 7 & nao \\
7 & D & 4 & 3 & 4 & nao \\
8 & D & 5 & 3 & 4 & nao \\
9 & D & 6 & 3 & 4 & nao \\
\bottomrule
\end{longtable}
\endgroup

\textbf{Resultado:}
\begin{itemize}[leftmargin=1.5em]
  \item valor na raiz continua 4 (mesma decisao final do Minimax);
  \item folhas podadas: 3 (as 3 ultimas folhas de $C$);
  \item ramos podados: no subarvore de $C$, apos avaliar a folha 1;
  \item nos nao explorados: 3 folhas (2, 0 e -1);
  \item nos visitados: 13 (4 internos + 9 folhas).
\end{itemize}

\textbf{Motivo matematico da poda:} em um no MIN, quando $\beta \le \alpha$, nenhum descendente restante pode melhorar a decisao de MAX na raiz.\\
\textbf{Observacao importante:} no registro do Alpha-Beta com poda, $C$ aparece com valor parcial 1, porque o no foi cortado apos a primeira folha. Isso nao contradiz o Minimax completo ($C=-1$), apenas indica avaliacao interrompida pelo limite $\beta \le \alpha$.

\subsection{Representacao visual da poda Alpha-Beta}
\begin{figure}[H]
\centering
\begin{forest}
for tree={
  draw,
  rounded corners,
  align=center,
  parent anchor=south,
  child anchor=north,
  l sep=10mm,
  s sep=7mm
}
[A\\MAX\\4
  [B\\MIN\\3
    [3]
    [5]
    [6]
    [9]
  ]
  [C\\MIN\\1*
    [1]
    [{\color{gray}2\\[-1pt]\tiny podado}, draw=gray, edge={gray,dashed}]
    [{\color{gray}0\\[-1pt]\tiny podado}, draw=gray, edge={gray,dashed}]
    [{\color{gray}-1\\[-1pt]\tiny podado}, draw=gray, edge={gray,dashed}]
  ]
  [D\\MIN\\4
    [7]
    [4]
    [5]
    [6]
  ]
]
\end{forest}
\caption{Poda Alpha-Beta na ordem original: em $C$, apenas a folha 1 e avaliada; as demais 3 folhas sao podadas}
\end{figure}

\subsection{Reordenacao para maximizar podas}
Reordenacao usada na implementacao:
\[
\text{raiz: } [D,B,C]
\]
e, dentro de cada no MIN, folhas em ordem crescente para provocar corte cedo.

\begin{itemize}[leftmargin=1.5em]
  \item Podas antes da reordenacao: 3 folhas podadas.
  \item Podas depois da reordenacao: 6 folhas podadas.
  \item Valor final na raiz permanece 4 (decisao correta preservada).
  \item A nova ordem melhora eficiencia porque encontra limites fortes cedo, ativando cortes antecipados.
  \item Relacao ordenacao x desempenho: quanto melhor a ordem de exploracao (melhores nos primeiro), maior tende a ser o ganho da poda Alpha-Beta.
\end{itemize}

\subsection{Minimax limitado (profundidade 2)}
Heuristicas de corte fornecidas:
\[
[4,7,2,5,6,1]
\]
Agrupando em 3 blocos MIN:
\[
\min(4,7)=4,\quad \min(2,5)=2,\quad \min(6,1)=1
\]
Logo, na raiz:
\[
\max(4,2,1)=4
\]
Decisao do limitado coincide com o valor da raiz no Minimax completo (4).
\textbf{Possiveis erros por heuristica:} se a estimativa local estiver enviesada, o Minimax limitado pode escolher um ramo diferente do Minimax completo, pois nao enxerga efeitos taticos que so aparecem em profundidades maiores.

\subsection{Experimento adicional: alterar 3 folhas}
Alteracao aplicada:
\[
[3,5,6] \rightarrow [8,8,8]
\]

\begin{itemize}[leftmargin=1.5em]
  \item Minimax modificado: raiz=8, $B=8$, $C=-1$, $D=4$.
  \item Alpha-Beta modificado: raiz=8, com 10 nos visitados e 6 folhas podadas.
  \item A decisao final muda: MAX passa de $D$ para $B$.
  \item Sensibilidade as utilidades: pequena alteracao em folhas do ramo B alterou o valor minimax do ramo de 3 para 8, mudando totalmente a decisao da raiz.
\end{itemize}

\subsection{Desenhos pos-experimento (Minimax e Alpha-Beta)}
\begin{figure}[H]
\centering
\resizebox{0.98\linewidth}{!}{
\begin{forest}
for tree={
  draw,
  rounded corners,
  align=center,
  parent anchor=south,
  child anchor=north,
  l sep=10mm,
  s sep=8mm
}
[A\\MAX\\8
  [B\\MIN\\8
    [8]
    [8]
    [8]
    [9]
  ]
  [C\\MIN\\-1
    [1]
    [2]
    [0]
    [-1]
  ]
  [D\\MIN\\4
    [7]
    [4]
    [5]
    [6]
  ]
]
\end{forest}
}
\caption{Arvore Minimax apos o experimento: MAX escolhe $B$ (valor 8)}
\end{figure}

\begin{figure}[H]
\centering
\resizebox{0.98\linewidth}{!}{
\begin{forest}
for tree={
  draw,
  rounded corners,
  align=center,
  parent anchor=south,
  child anchor=north,
  l sep=10mm,
  s sep=7mm
}
[A\\MAX\\8
  [B\\MIN\\8
    [8]
    [8]
    [8]
    [9]
  ]
  [C\\MIN\\1*
    [1]
    [{\color{gray}2\\[-1pt]\tiny podado}, draw=gray, edge={gray,dashed}]
    [{\color{gray}0\\[-1pt]\tiny podado}, draw=gray, edge={gray,dashed}]
    [{\color{gray}-1\\[-1pt]\tiny podado}, draw=gray, edge={gray,dashed}]
  ]
  [D\\MIN\\7*
    [7]
    [{\color{gray}4\\[-1pt]\tiny podado}, draw=gray, edge={gray,dashed}]
    [{\color{gray}5\\[-1pt]\tiny podado}, draw=gray, edge={gray,dashed}]
    [{\color{gray}6\\[-1pt]\tiny podado}, draw=gray, edge={gray,dashed}]
  ]
]
\end{forest}
}
\caption{Alpha-Beta pos-experimento (ordem original): 6 podas (3 em $C$ e 3 em $D$)}
\end{figure}

\subsection{Comparacao final (Q5)}
\begin{table}[H]
\centering
\begin{tabular}{l c c}
\toprule
Cenario & Minimax (nos) & Alpha-Beta (nos / podas) \\
\midrule
Original & 16 & 13 / 3 folhas \\
Reordenado & - & 10 / 6 folhas \\
Modificado (3 folhas) & 16 & 10 / 6 folhas \\
\bottomrule
\end{tabular}
\caption{Comparacao quantitativa das execucoes}
\end{table}

\begin{table}[H]
\centering
\begingroup
\footnotesize
\setlength{\tabcolsep}{4pt}
\renewcommand{\arraystretch}{1.18}
\begin{tabular}{>{\raggedright\arraybackslash}p{3.6cm} >{\raggedright\arraybackslash}p{4.5cm} >{\raggedright\arraybackslash}p{4.5cm}}
\toprule
\textbf{Criterio} & \textbf{Minimax} & \textbf{Alpha-Beta} \\
\midrule
Nos explorados (nesta prova) & 16 nos no cenario original e 16 no modificado & 13 nos no original; 10 com reordenacao e no modificado \\
Quantidade de podas & 0 (nao realiza poda) & 3 podas no original; 6 apos reordenacao \\
Custo computacional (teorico) & $O(b^m)$ & pior caso $O(b^m)$; melhor caso com boa ordenacao $O(b^{m/2})$ \\
Consumo de memoria (teorico) & $O(bm)$ (busca em profundidade) & $O(bm)$, mesma ordem assintotica \\
Impacto da ordenacao dos filhos & baixo (nao altera exploracao) & alto (ordem melhor aumenta as podas) \\
Qualidade da decisao & exata (adversario perfeito) & exata e igual ao Minimax quando a poda e correta \\
\bottomrule
\end{tabular}
\endgroup
\caption{Comparacao Minimax vs Alpha-Beta com criterios teoricos e evidencias da execucao}
\end{table}

\textbf{Sintese:}
\begin{itemize}[leftmargin=1.5em]
  \item custo computacional: Alpha-Beta $\le$ Minimax (mesma decisao, menos nos);
  \item consumo de memoria: mesma ordem assintotica (DFS na arvore), diferindo pouco na pratica aqui;
  \item impacto da ordenacao: muito alto no Alpha-Beta (podas sobem de 3 para 6);
  \item qualidade da decisao: igual ao Minimax quando os cortes sao corretos.
\end{itemize}
